\section{Building the Time-Evolution MPO}
\label{sec:mporep}

The transverse contraction framework introduced in Section~\ref{sec:transverse} requires the short-time evolution operator $U(\delta t)=e^{-iH\delta t}$ to be represented as a translation-invariant \ac{MPO}. This ensures that every column of the rotated space-time network defines the same transfer matrix, effectively reducing the spatial contraction (which encodes the dynamics) to an eigenvalue problem. If the local tensors depended on the site, each column would instead define a different operator and this simple reduction would no longer apply.

For \ac{NN} Hamiltonians, such an operator is usually constructed through a Trotter--Suzuki decomposition into layers of two-site gates, as summarized in Appendix~\ref{app:mpo:standard}. A term coupling sites $i$ and $i+2$ does not fit into this bond pattern, and the standard extensions do not preserve a one-site translation-invariant \ac{MPO}. We therefore represent the Hamiltonian as a \acf{FSM} and exponentiate this representation directly.

\subsection{The Hamiltonian as a \texorpdfstring{\ac{FSM}}{FSM}}
\label{sec:mporep:fsm}

In the \ac{FSM} representation, the Hamiltonian is encoded as an \ac{MPO} whose local tensor has a block-upper-triangular form. This tensor is composed of two identity blocks, which propagate the initial and final states; an on-site block $D$; a block row $C$, which initiates interactions; a block column $B$, which terminates them; and a memory block $A$, which propagates unfinished interactions between their endpoints. The virtual index therefore records whether an interaction has not yet begun, is being propagated, or has terminated. For the ANNNI-type model, the local tensor has bond dimension $D_W=5$ and takes the form
\begin{align}
  W_H & =
  \begin{pmatrix}
    \mathbbm{1} & C & D \\
    0 & A & B \\
    0 & 0 & \mathbbm{1}
  \end{pmatrix}=
  \begin{pmatrix}
    \mathbbm{1} & -\sigma^z & -p\lambda\,\sigma^x & -p\,\sigma^z & -\lambda\,\sigma^x \\
    0 & 0 & 0 & 0 & \sigma^z \\
    0 & 0 & 0 & 0 & \sigma^x \\
    0 & \mathbbm{1} & 0 & 0 & 0 \\
    0 & 0 & 0 & 0 & \mathbbm{1}
  \end{pmatrix}.
  \label{eq:Wannni}
\end{align}

Figure~\ref{fig:fsm} represents Equation~\eqref{eq:Wannni} as the allowed paths through the five virtual states. Each node denotes a value of the virtual index, while an arrow from state $a$ to state $b$, labelled by an operator $O$, represents the tensor entry $(W_H)_{ab}=O$. The machine scans the chain from left to right, beginning in state $1$ and ending in state $5$.
\begin{figure}[htbp]
  \centering
  \includegraphics[width=0.8\linewidth]{imgs/fsm.pdf}%
  \caption{\ac{FSM} representation of the Hamiltonian \ac{MPO}. The nodes are the five virtual states of Equation~\eqref{eq:Wannni}, and each arrow is labelled by the physical operator applied during the corresponding transition. The machine starts in state $1$ and closes a term on reaching state $5$, drawn with a double circle; every path between them generates one term of the Hamiltonian. Arrows are coloured by the term they build: the transverse field in teal, the two nearest-neighbour couplings in blue ($\sigma^z\sigma^z$) and orange ($\sigma^x\sigma^x$), and the \ac{NNN} coupling in purple. The grey loops are the identities applied before the first term starts and after the last one closes. State $2$ is shared by the blue and purple paths because both terminate on a $\sigma^z$, which is what keeps the bond dimension at $D_W=5$.}
  \label{fig:fsm}
\end{figure}
The loop $1\to1$ applies identities until a Hamiltonian term is initiated. A direct transition $1\to5$ applies the on-site field and completes the term at the same site. For an \ac{NN} interaction, the machine instead follows either $1\to2\to5$ or $1\to3\to5$. The first transition applies the operator at site $i$ and stores in the virtual state whether a $\sigma^z$ or $\sigma^x$ endpoint is pending; the second applies the matching operator at site $i+1$ and completes the interaction. The \ac{NNN} term follows the path $1\to4\to2\to5$: the transition to state $4$ applies the first $\sigma^z$ at site $i$, the transition $4\to2$ inserts an identity at site $i+1$, and the transition $2\to5$ applies the final $\sigma^z$ at site $i+2$. Once state $5$ is reached, the loop $5\to5$ applies identities along the remainder of the chain. The zero entries of $W_H$ correspond to absent arrows, so no other transitions are allowed.

This directional construction makes the tensor asymmetric in its virtual indices. After the $90^\circ$ rotation, the corresponding spatial transfer matrix is consequently non-Hermitian and non-normal. The left and right temporal boundary states are therefore independent and must be computed separately. Section~\ref{sec:nonhermitian} introduces the resulting notation.

\subsection{Exponentiating the \texorpdfstring{\ac{FSM}}{FSM}}
\label{sec:mporep:exp}

We next construct an approximation $W(\tau)\approx e^{\tau H}$, with $\tau=-i\delta t$, from the blocks of $W_H$. This operation is not equivalent to exponentiating the local block matrix: powers of an upper-triangular matrix mix its blocks and generate products such as $CD$, $BC$, and $AD$, whose coefficients must reproduce the Taylor expansion to the required order. Directly truncating the Taylor series of the many-body operator is also unsuitable because $H^n$ contains $\mathcal{O}(N^n)$ local contributions, so the resulting approximation is not size-extensive~\cite{vandamme2024}. The constructions used here instead organize the expansion by the spatial overlap of the local terms. Contributions with disjoint support commute and factorize, allowing them to be retained to all orders, while only connected clusters of overlapping terms need to be reproduced up to a prescribed size. Appendix~\ref{app:mpo} gives the detailed construction.

We consider two families of approximations built from these blocks. The operators $W^{\mathrm{I}}$ and $W^{\mathrm{II}}$ introduced in Ref.~\cite{zaletel2015} have virtual bond dimension $1+\chi$, where $\chi=3$ is the number of memory channels in the present model. Although $W^{\mathrm{II}}$ is effectively second order for a strictly \ac{NN} Hamiltonian, it omits connected clusters that arise once the \ac{NNN} interaction is included and is therefore only first order for our model. We instead use the higher-order construction of Ref.~\cite{vandamme2024}, whose order-$n$ approximation includes the connected clusters required to produce a one-step error $\mathcal{O}(\delta t^{n+1})$. The order-1 approximation coincides with $W^{\mathrm{II}}$, while the order-2 approximation used throughout this work is denoted VD2.

VD2 increases the virtual bond dimension from $1+\chi$ to $1+\chi+\chi^2$, which gives $13$ rather than $4$ for this model. After rotating the network, this bond index becomes the physical index of each temporal site, so the \ac{tMPS} has a larger local dimension than it would with $W^{\mathrm{II}}$. In return, second-order accuracy allows us to use the comparatively large time step $\delta t=0.05$ while controlling the discretization error. This trade-off is important because reducing $\delta t$ increases the number of temporal sites required to reach a fixed physical time, lengthening the temporal chain and introducing additional truncation steps. The second-order operator therefore makes the long-time calculations of Section~\ref{sec:results} computationally accessible despite its larger local dimension. Appendix~\ref{app:mpo:comparison} compares the two constructions and verifies the first-order convergence of $W^{\mathrm{II}}$ for the ANNNI-type model.

\subsection{The non-Hermitian transfer matrix}
\label{sec:nonhermitian}

Rotating the space-time network turns the evolution operator into the spatial transfer matrix $\mathcal{E}$. Because it inherits the virtual asymmetry of the \ac{FSM}, this matrix is non-Hermitian and non-normal, $\mathcal{E}\mathcal{E}^\dagger\neq\mathcal{E}^\dagger\mathcal{E}$. These properties were used in Section~\ref{sec:transverse}; here we define the notation needed for their numerical treatment.

The right and left eigenvectors satisfy $\mathcal{E}\ket{R_i}=\mu_i\ket{R_i}$ and $\bra{L_i}\mathcal{E}=\mu_i\bra{L_i}$, respectively, and are generally unrelated. Neither eigenbasis is orthonormal by itself. Instead, the two sets are paired through the biorthogonality relation $\braket{L_i}{R_j}\propto\delta_{ij}$, where the relevant pairing is the bilinear contraction defined by the tensor network rather than the usual Hermitian inner product. We order the generally complex eigenvalues $\mu_i$ by modulus and denote the dominant one by $\mu_0$. This convention is used throughout Sections~\ref{sec:numerics} and~\ref{sec:results}.
